% =========================================================
% Ordered Pairs and Relations
% =========================================================

\subsection{Ordered Pairs and Relations}

% ---------------------------------------------------------
% TOOLKIT
% ---------------------------------------------------------

\begin{remark*}[Geometric picture]
A binary relation from $A$ to $B$ can be pictured as directed edges from the
$A$-side to the $B$-side of a bipartite diagram. The converse reverses every
edge, and composition follows two edges in series through an intermediate set.
\end{remark*}

\begin{definitionbox}{Definition (Ordered Pair)}
\begin{definition}[Ordered Pair]\label{def:ordered-pair}
Let $a$ and $b$ be sets. The \emph{ordered pair} $(a,b)$ is defined (Kuratowski)
as
\[
(a,b) \;:=\; \{\{a\}, \{a,b\}\}.
\]
\end{definition}
\end{definitionbox}

\begin{dependencies}
\begin{itemize}
  \item \hyperref[ax:pairing]{Axiom of Pairing}
\end{itemize}
\end{dependencies}

\begin{theorembox}{Theorem (Uniqueness of Ordered Pairs)}
\begin{theorem}[Uniqueness of Ordered Pairs]\label{thm:ordered-pair-unique}
For any sets $a,b,c,d$,
\[
(a,b) = (c,d)
\;\;\Longleftrightarrow\;\;
(a = c \land b = d).
\]
\hyperref[prf:ordered-pair-unique]{Go to proof.}
\end{theorem}
\end{theorembox}

\begin{remark*}[Standard quantified statement]
$\forall a, b, c, d:\; (a,b) = (c,d) \;\leftrightarrow\; (a = c \wedge b = d)$.
\end{remark*}
\begin{dependencies}
\begin{itemize}
  \item \hyperref[def:ordered-pair]{Ordered Pair}
  \item \hyperref[ax:pairing]{Axiom of Pairing}
  \item \hyperref[ax:extensionality]{Axiom of Extensionality}
\end{itemize}
\end{dependencies}

\begin{remark*}[Kuratowski encoding]
The purpose of the Kuratowski definition is purely to guarantee the uniqueness
theorem above. Once $(a,b) = (c,d) \iff a=c \land b=d$ is established,
we may treat ordered pairs as a primitive with this property and forget the
encoding.
\end{remark*}

\begin{definitionbox}{Definition (Cartesian Product --- as foundation for relations)}
\begin{definition}[Cartesian Product --- as foundation for relations]\label{def:cartesian-product-rel}
The \emph{Cartesian product} of sets $A$ and $B$ is:
\[
A \times B
\;:=\;
\{\, (a,b) \mid a \in A \text{ and } b \in B \,\}.
\]
\end{definition}
\end{definitionbox}

\begin{dependencies}
\begin{itemize}
  \item \hyperref[def:ordered-pair]{Ordered Pair}
  \item \hyperref[ax:pairing]{Axiom of Pairing}
  \item \hyperref[ax:power-set]{Axiom of Power Set}
  \item \hyperref[ax:separation]{Axiom Schema of Separation}
\end{itemize}
\end{dependencies}

\begin{definitionbox}{Definition (Relation)}
\begin{definition}[Relation]\label{def:relation}
A \emph{relation} from $A$ to $B$ is any subset $R \subseteq A \times B$.

If $(a,b) \in R$ we write $a \, R \, b$ and say ``$a$ is related to $b$.''
When $A = B$ we say $R$ is a \emph{relation on $A$}.
\end{definition}
\end{definitionbox}


\begin{dependencies}
\begin{itemize}
  \item \hyperref[def:cartesian-product-rel]{Cartesian Product}
  \item \hyperref[def:subset]{Subset}
\end{itemize}
\end{dependencies}

\begin{remark*}[Relations as structured sets]
Relations introduce no new foundational objects: they are simply sets of
ordered pairs, constructed using operations already available. Properties of
relations can therefore be studied with ordinary set-theoretic reasoning.
\end{remark*}

\begin{definitionbox}{Definition (Domain of a Relation)}
\begin{definition}[Domain of a Relation]\label{def:relation-domain}
Let $R\subseteq A\times B$ be a relation. The \emph{domain} of $R$ is
\[
\operatorname{Dom}(R)
\;:=\;
\{\,a\in A\mid \exists b\in B,\ (a,b)\in R\,\}.
\]
\end{definition}
\end{definitionbox}

\begin{dependencies}
\begin{itemize}
  \item \hyperref[def:relation]{Relation}
  \item \hyperref[ax:separation]{Axiom Schema of Separation}
\end{itemize}
\end{dependencies}

\begin{definitionbox}{Definition (Range of a Relation)}
\begin{definition}[Range of a Relation]\label{def:relation-range}
Let $R\subseteq A\times B$ be a relation. The \emph{range} of $R$ is
\[
\operatorname{Ran}(R)
\;:=\;
\{\,b\in B\mid \exists a\in A,\ (a,b)\in R\,\}.
\]
Tarski also calls this class the \emph{counter-domain} or \emph{converse domain}.
\end{definition}
\end{definitionbox}

\begin{dependencies}
\begin{itemize}
  \item \hyperref[def:relation]{Relation}
  \item \hyperref[ax:separation]{Axiom Schema of Separation}
\end{itemize}
\end{dependencies}

\begin{remark*}[Domain and range]
The domain records the first coordinates that occur in a relation; the range
records the second coordinates that occur.
\end{remark*}

\begin{definitionbox}{Definition (Converse Relation)}
\begin{definition}[Converse Relation]\label{def:converse-relation}
If $R\subseteq A\times B$, the \emph{converse} of $R$ is the relation
$R^{-1}\subseteq B\times A$ defined by
\[
(b,a)\in R^{-1}
\;\Longleftrightarrow\;
(a,b)\in R.
\]
Equivalently, $b\,R^{-1}\,a$ iff $aRb$.
\end{definition}
\end{definitionbox}

\begin{dependencies}
\begin{itemize}
  \item \hyperref[def:relation]{Relation}
  \item \hyperref[def:ordered-pair]{Ordered Pair}
\end{itemize}
\end{dependencies}

\begin{definitionbox}{Definition (Identity Relation)}
\begin{definition}[Identity Relation]\label{def:relation-identity}
For a set $A$, the \emph{identity relation} on $A$ is
\[
\Delta_A
\;:=\;
\{\, (a,a)\mid a\in A\,\}.
\]
\end{definition}
\end{definitionbox}

\begin{dependencies}
\begin{itemize}
  \item \hyperref[def:relation]{Relation}
  \item \hyperref[def:cartesian-product-rel]{Cartesian Product}
\end{itemize}
\end{dependencies}

\begin{definitionbox}{Definition (Composition of Relations)}
\begin{definition}[Composition of Relations]\label{def:relation-composition}
Let $R\subseteq A\times B$ and $S\subseteq B\times C$. The
\emph{composition} $S\circ R\subseteq A\times C$ is defined by
\[
(a,c)\in S\circ R
\;\Longleftrightarrow\;
\exists b\in B,\; (a,b)\in R \text{ and } (b,c)\in S.
\]
\end{definition}
\end{definitionbox}

\begin{dependencies}
\begin{itemize}
  \item \hyperref[def:relation]{Relation}
  \item \hyperref[def:cartesian-product-rel]{Cartesian Product}
\end{itemize}
\end{dependencies}

\begin{theorembox}{Theorem (Associativity of Relation Composition)}
\begin{theorem}[Associativity of Relation Composition]\label{thm:relation-composition-assoc}
Let $R\subseteq A\times B$, $S\subseteq B\times C$, and
$T\subseteq C\times D$. Then
\[
T\circ(S\circ R)=(T\circ S)\circ R.
\]
\hyperref[prf:relation-composition-assoc]{Go to proof.}
\end{theorem}
\end{theorembox}

\begin{remark*}[Standard quantified statement]
For all $a\in A$ and $d\in D$, $(a,d)\in T\circ(S\circ R)$ iff there are
$b\in B$ and $c\in C$ with $(a,b)\in R$, $(b,c)\in S$, and $(c,d)\in T$,
which is exactly the membership condition for $(a,d)\in (T\circ S)\circ R$.
\end{remark*}

\begin{dependencies}
\begin{itemize}
  \item \hyperref[def:relation-composition]{Composition of Relations}
  \item \hyperref[def:relation-equality]{Relation Equality}
\end{itemize}
\end{dependencies}

\begin{theorembox}{Theorem (Identity Laws for Relation Composition)}
\begin{theorem}[Identity Laws for Relation Composition]\label{thm:relation-composition-id}
Let $R\subseteq A\times B$. Then
\[
R\circ\Delta_A=R
\quad\text{and}\quad
\Delta_B\circ R=R.
\]
\hyperref[prf:relation-composition-id]{Go to proof.}
\end{theorem}
\end{theorembox}

\begin{dependencies}
\begin{itemize}
  \item \hyperref[def:relation-identity]{Identity Relation}
  \item \hyperref[def:relation-composition]{Composition of Relations}
\end{itemize}
\end{dependencies}

\begin{definitionbox}{Definition (Null Relation)}
\begin{definition}[Null Relation]\label{def:null-relation}
For sets $A$ and $B$, the \emph{null relation} from $A$ to $B$ is
$\varnothing\subseteq A\times B$. No element of $A$ is related to any element
of $B$.
\end{definition}
\end{definitionbox}

\begin{dependencies}
\begin{itemize}
  \item \hyperref[def:relation]{Relation}
  \item \hyperref[def:empty-set]{Empty Set}
  \item \hyperref[def:cartesian-product-rel]{Cartesian Product}
\end{itemize}
\end{dependencies}

\begin{definitionbox}{Definition (Universal Relation)}
\begin{definition}[Universal Relation]\label{def:universal-relation}
For sets $A$ and $B$, the \emph{universal relation} from $A$ to $B$ is
$A\times B$ itself. Every element of $A$ is related to every element of $B$.
\end{definition}
\end{definitionbox}

\begin{dependencies}
\begin{itemize}
  \item \hyperref[def:relation]{Relation}
  \item \hyperref[def:cartesian-product-rel]{Cartesian Product}
\end{itemize}
\end{dependencies}

\begin{remark*}[Extreme relations]
The null relation and universal relation are the two extreme binary relations
from $A$ to $B$.
\end{remark*}

\begin{definitionbox}{Definition (Relation Inclusion)}
\begin{definition}[Relation Inclusion]\label{def:relation-inclusion}
Let $R,S\subseteq A\times B$. Relation inclusion is ordinary set inclusion:
\[
R\subseteq S
\;\Longleftrightarrow\;
\forall a\in A\,\forall b\in B,\ ((a,b)\in R\Rightarrow (a,b)\in S),
\]
\end{definition}
\end{definitionbox}

\begin{dependencies}
\begin{itemize}
  \item \hyperref[def:relation]{Relation}
  \item \hyperref[def:subset]{Subset}
\end{itemize}
\end{dependencies}

\begin{definitionbox}{Definition (Relation Equality)}
\begin{definition}[Relation Equality]\label{def:relation-equality}
Let $R,S\subseteq A\times B$. Relation equality is ordinary set equality:
\[
R=S
\;\Longleftrightarrow\;
R\subseteq S \text{ and } S\subseteq R.
\]
\end{definition}
\end{definitionbox}

\begin{dependencies}
\begin{itemize}
  \item \hyperref[def:relation]{Relation}
  \item \hyperref[def:relation-inclusion]{Relation Inclusion}
  \item \hyperref[def:set-equality]{Set Equality}
\end{itemize}
\end{dependencies}

\begin{definitionbox}{Definition (Union of Relations)}
\begin{definition}[Union of Relations]\label{def:relation-union}
Let $R,S\subseteq A\times B$. The \emph{union} $R\cup S$ is the relation with
\[
(a,b)\in R\cup S
\;\Longleftrightarrow\;
(a,b)\in R \text{ or } (a,b)\in S.
\]
\end{definition}
\end{definitionbox}

\begin{dependencies}
\begin{itemize}
  \item \hyperref[def:relation]{Relation}
  \item \hyperref[def:union]{Union}
\end{itemize}
\end{dependencies}

\begin{definitionbox}{Definition (Intersection of Relations)}
\begin{definition}[Intersection of Relations]\label{def:relation-intersection}
Let $R,S\subseteq A\times B$. The \emph{intersection} $R\cap S$ is the relation
with
\[
(a,b)\in R\cap S
\;\Longleftrightarrow\;
(a,b)\in R \text{ and } (a,b)\in S.
\]
\end{definition}
\end{definitionbox}

\begin{dependencies}
\begin{itemize}
  \item \hyperref[def:relation]{Relation}
  \item \hyperref[def:intersection]{Intersection}
\end{itemize}
\end{dependencies}

\begin{definitionbox}{Definition (Difference of Relations)}
\begin{definition}[Difference of Relations]\label{def:relation-difference}
Let $R,S\subseteq A\times B$. The \emph{difference} $R\setminus S$ is the
relation with
\[
(a,b)\in R\setminus S
\;\Longleftrightarrow\;
(a,b)\in R \text{ and } (a,b)\notin S.
\]
\end{definition}
\end{definitionbox}

\begin{dependencies}
\begin{itemize}
  \item \hyperref[def:relation]{Relation}
  \item \hyperref[def:set-difference]{Set Difference}
\end{itemize}
\end{dependencies}

\begin{definitionbox}{Definition (Complement of a Relation)}
\begin{definition}[Complement of a Relation]\label{def:relation-complement}
Let $R\subseteq A\times B$. The \emph{complement} of $R$ relative to
$A\times B$ is
\[
(A\times B)\setminus R.
\]
\end{definition}
\end{definitionbox}

\begin{dependencies}
\begin{itemize}
  \item \hyperref[def:relation]{Relation}
  \item \hyperref[def:set-difference]{Set Difference}
\end{itemize}
\end{dependencies}

\begin{remark*}[Relation algebra operations]
Relation inclusion, equality, union, intersection, difference, and complement
are inherited from the corresponding set operations on subsets of $A\times B$.
\end{remark*}

\begin{theorembox}{Theorem (Converse Laws for Relation Operations)}
\begin{theorem}[Converse Laws for Relation Operations]\label{thm:relation-converse-laws}
Let $R,S\subseteq A\times B$, let $T\subseteq B\times C$, and take complements
relative to the indicated Cartesian products. Then
\[
(R^{-1})^{-1}=R,\qquad
(T\circ R)^{-1}=R^{-1}\circ T^{-1},
\]
\[
(R\cup S)^{-1}=R^{-1}\cup S^{-1},\qquad
(R\cap S)^{-1}=R^{-1}\cap S^{-1},
\]
\[
(R\setminus S)^{-1}=R^{-1}\setminus S^{-1},\qquad
((A\times B)\setminus R)^{-1}=(B\times A)\setminus R^{-1}.
\]
\hyperref[prf:relation-converse-laws]{Go to proof.}
\end{theorem}
\end{theorembox}

\begin{dependencies}
\begin{itemize}
  \item \hyperref[def:converse-relation]{Converse Relation}
  \item \hyperref[def:relation-composition]{Composition of Relations}
  \item \hyperref[def:relation-union]{Union of Relations}
  \item \hyperref[def:relation-intersection]{Intersection of Relations}
  \item \hyperref[def:relation-difference]{Difference of Relations}
  \item \hyperref[def:relation-complement]{Complement of a Relation}
\end{itemize}
\end{dependencies}

\begin{theorembox}{Theorem (Relation Composition and Boolean Operations)}
\begin{theorem}[Relation Composition and Boolean Operations]\label{thm:relation-composition-boolean}
Let $P,Q\subseteq A\times B$, $R\subseteq B\times C$, and
$S,T\subseteq B\times C$. Then
\[
R\circ(P\cup Q)=(R\circ P)\cup(R\circ Q),
\qquad
(S\cup T)\circ P=(S\circ P)\cup(T\circ P).
\]
Also,
\[
R\circ(P\cap Q)\subseteq(R\circ P)\cap(R\circ Q),
\qquad
(S\cap T)\circ P\subseteq(S\circ P)\cap(T\circ P).
\]
The analogous formulas for difference and complement hold only as inclusions
or under extra hypotheses; they are not general distributive laws.
\hyperref[prf:relation-composition-boolean]{Go to proof.}
\end{theorem}
\end{theorembox}

\begin{dependencies}
\begin{itemize}
  \item \hyperref[def:relation-composition]{Composition of Relations}
  \item \hyperref[def:relation-union]{Union of Relations}
  \item \hyperref[def:relation-intersection]{Intersection of Relations}
\end{itemize}
\end{dependencies}

\begin{remark*}[Warning on relation composition]
Relation composition distributes over unions on both sides. It does not
generally distribute over intersections, differences, or complements. The
reason is existential: a pair can lie in both composed relations using two
different middle witnesses, even when no single witness satisfies both
conditions at once.
\end{remark*}

\begin{remark*}[Relation algebra interaction table]
\mbox{}\par\smallskip
\small
\begin{tabularx}{\textwidth}{l X X}
\toprule
\textbf{Operation} & \textbf{Converse} & \textbf{Composition} \\
\midrule
Union
& $(R\cup S)^{-1}=R^{-1}\cup S^{-1}$
& Distributes exactly on both sides \\
Intersection
& $(R\cap S)^{-1}=R^{-1}\cap S^{-1}$
& Gives $\subseteq$ in general; equality needs extra hypotheses \\
Difference
& $(R\setminus S)^{-1}=R^{-1}\setminus S^{-1}$
& No general distributive equality \\
Complement
& $((A\times B)\setminus R)^{-1}=(B\times A)\setminus R^{-1}$
& No general distributive equality \\
Composition
& $(T\circ R)^{-1}=R^{-1}\circ T^{-1}$
& Associative with identity relation $\Delta_A$ \\
\bottomrule
\end{tabularx}
\end{remark*}

\begin{definitionbox}{Definition (Transitive Closure)}
\begin{definition}[Transitive Closure]\label{def:relation-transitive-closure}
Let $R$ be a relation on $A$. The \emph{transitive closure} of $R$ is the
smallest transitive relation on $A$ containing $R$.
\end{definition}
\end{definitionbox}

\begin{dependencies}
\begin{itemize}
  \item \hyperref[def:relation]{Relation}
  \item \hyperref[def:transitive]{Transitive Relation}
  \item \hyperref[def:relation-inclusion]{Relation Inclusion}
\end{itemize}
\end{dependencies}

\begin{definitionbox}{Definition (Equivalence Closure)}
\begin{definition}[Equivalence Closure]\label{def:relation-equivalence-closure}
Let $R$ be a relation on $A$. The \emph{equivalence closure} of $R$ is the
smallest equivalence relation on $A$ containing $R$.
\end{definition}
\end{definitionbox}

\begin{dependencies}
\begin{itemize}
  \item \hyperref[def:relation]{Relation}
  \item \hyperref[def:equivalence-rel]{Equivalence Relation}
  \item \hyperref[def:relation-inclusion]{Relation Inclusion}
\end{itemize}
\end{dependencies}

\begin{definitionbox}{Definition (One-to-Many Relation)}
\begin{definition}[One-to-Many Relation]\label{def:one-to-many-relation}
Let $R\subseteq A\times B$.
$R$ is \emph{one-to-many} if some input is related to two distinct outputs:
\[
\exists a\in A\,\exists b_1,b_2\in B,\quad
(a,b_1)\in R\land (a,b_2)\in R\land b_1\neq b_2.
\]
\end{definition}
\end{definitionbox}

\begin{dependencies}
\begin{itemize}
  \item \hyperref[def:relation]{Relation}
\end{itemize}
\end{dependencies}

\begin{definitionbox}{Definition (Many-to-One Relation)}
\begin{definition}[Many-to-One Relation]\label{def:many-to-one-relation}
Let $R\subseteq A\times B$.
$R$ is \emph{many-to-one} if two distinct inputs are related to a common output:
\[
\exists a_1,a_2\in A\,\exists b\in B,\quad
(a_1,b)\in R\land (a_2,b)\in R\land a_1\neq a_2.
\]
\end{definition}
\end{definitionbox}

\begin{dependencies}
\begin{itemize}
  \item \hyperref[def:relation]{Relation}
\end{itemize}
\end{dependencies}

\begin{definitionbox}{Definition (Many-to-Many Relation)}
\begin{definition}[Many-to-Many Relation]\label{def:many-to-many-relation}
Let $R\subseteq A\times B$.
$R$ is \emph{many-to-many} if it is one-to-many and many-to-one.
\end{definition}
\end{definitionbox}

\begin{dependencies}
\begin{itemize}
  \item \hyperref[def:one-to-many-relation]{One-to-Many Relation}
  \item \hyperref[def:many-to-one-relation]{Many-to-One Relation}
\end{itemize}
\end{dependencies}

\begin{remark*}[Relation cardinality patterns]
The one-to-many, many-to-one, and many-to-many labels classify how many inputs
and outputs can be connected by a relation.
\end{remark*}

\begin{remark*}[Why one-to-many is not functional]
A one-to-many relation fails the determinacy condition needed for a function:
one input admits at least two possible outputs. A many-to-one relation may
still be a function, because different inputs are allowed to share an output.
\end{remark*}

\begin{definitionbox}{Definition (Many-Place Relation)}
\begin{definition}[Many-Place Relation]\label{def:many-place-relation}
A \emph{ternary relation} among sets $A$, $B$, and $C$ is a subset of
$A\times B\times C$. More generally, an $n$-place relation on
$A_1,\ldots,A_n$ is a subset of
\[
A_1\times\cdots\times A_n.
\]
For example, betweenness in geometry is naturally a three-place relation, and
the arithmetic statement $x+y=z$ is naturally a three-place relation among
$x$, $y$, and $z$.
\end{definition}
\end{definitionbox}

\begin{dependencies}
\begin{itemize}
  \item \hyperref[def:relation]{Relation}
  \item \hyperref[def:cartesian-product-rel]{Cartesian Product}
\end{itemize}
\end{dependencies}

\begin{remark*}[Tarski-style hierarchy]
The hierarchy used in the next chapter is:
\[
\begin{aligned}
\text{relation}
&\longrightarrow
\text{functional relation}
\longrightarrow
\text{function} \\
&\longrightarrow
\text{one-one or onto function}
\longrightarrow
\text{bijection}.
\end{aligned}
\]
The first step adds determinacy; the second adds totality; the later steps
control whether outputs are shared or missed.
\end{remark*}
